Photocurrent spectroscopy is a rather unconventional type of ultrafast pump-probe spectroscopy technique where instead of analyzing a probe beam containing information about a system one gathers information by analyzing a photocurrent. This photocurrent is either fully realized by laser excitation (photo-current spectroscopy), or an externally applied gating current is altered by modifying the materials conductivity by laser excitation (Photo-conductivity spectroscopy). There are many different variations of this technique, which has existed since the $1982$, but only became popular around $2012$ due to ..... The version we use is a far-field implementation of photocurrent cross-correlation spectroscopy.\\
\newline
In photocurrent cross-correlation spectroscopy one uses two interferometrically seperated light pulses to successively excite a sample. 
